\section{Safety Analyses}

Safety analyses will be performed in the Safety population using descriptive statistics
only; no formal hypothesis testing is pre-specified for safety endpoints.

\subsection{Adverse Events}

\begin{itemize}[leftmargin=1.6em, itemsep=2pt]
  \item Adverse events (AEs) are coded using MedDRA v27.0.
  \item Incidence of Treatment-Emergent AEs (TEAEs) summarised by System Organ Class
        (SOC) and Preferred Term (PT).
  \item Participants counted once per PT at the maximum severity and strongest relationship
        to study drug.
  \item Separate summaries for: all TEAEs, Grade $\geq$3 TEAEs, SAEs, TEAEs leading to
        discontinuation, AEs of Special Interest (AESI: cardiovascular, hepatic).
\end{itemize}

\subsection{Clinical Laboratory Parameters}

Continuous lab parameters (haematology, biochemistry, urinalysis): descriptive statistics
(n, mean, SD, median, Q1/Q3, min, max) at each visit. Shift tables (normal/abnormal at
baseline vs.\ post-baseline) for key parameters. Notable abnormalities flagged using
CTCAE v5.0 grading and sponsor-defined reference ranges.

\subsection{Vital Signs and ECG}

Vital signs (blood pressure, heart rate, respiratory rate, SpO\textsubscript{2}) and
12-lead ECG parameters summarised descriptively. QTcF prolongation: proportion of
participants with QTcF $>$450, $>$480, $>$500\,ms and change from baseline $>$30, $>$60\,ms.

\begin{tcolorbox}[sapwarn, title={\faIcon{exclamation-triangle}\ Safety Review Committee Note}]
  An independent Safety Review Committee (SRC) will review unblinded safety data at
  Weeks 12 and 26. SRC findings may trigger an ad hoc safety analysis per the SRC
  Charter (Doc.\ ARI-301-SRC-001). Such analyses are not pre-specified in this SAP
  and will be reported separately.
\end{tcolorbox}

\clearpage

% ============================================================
%  SECTION 7 — MISSING DATA HANDLING
% ============================================================
